%% ma: L12-B06-0054
%% lop: 12
%% chuong: 2
%% bai: 6
%% dang: 12.6.20
%% loai: TN4
%% muc: TH
%% dap_an: "B"
%% nguon: "CĐ2.1.pdf giáo viên gửi 27/10/2026 (M07, Lê Minh Khôi – Nguyễn Anh Khôi, Chương II Chủ đề 1 Vectơ trong không gian), Câu 14"
%% kien_thuc: "Bài 6 (trọng tâm tam giác, trung điểm)"
%% trang_thai: cho-duyet
%% ly_do: "Dạng mới 12.6.20 đề xuất, chờ giáo viên duyệt"
\begin{ex}
Cho tứ diện $ABCD$ có $M$ và $N$ lần lượt là trung điểm của $AB$ và $CD$, $G$ là trọng tâm của tam giác $BCD$. Mệnh đề nào sau đây đúng?
\choice{$\overrightarrow{MN}=\dfrac13\cdot\left(\overrightarrow{AB}+\overrightarrow{CD}\right)$}{\True $\overrightarrow{AB}+\overrightarrow{AC}+\overrightarrow{AD}=3\cdot\overrightarrow{AG}$}{$\overrightarrow{GD}-2\cdot\overrightarrow{GN}=\vec 0$}{$\overrightarrow{AM}+\overrightarrow{AN}=\dfrac23\cdot\overrightarrow{AG}$}
\loigiai{$G$ là trọng tâm tam giác $BCD$ nên $\overrightarrow{GB}+\overrightarrow{GC}+\overrightarrow{GD}=\vec 0$, suy ra $\overrightarrow{AB}+\overrightarrow{AC}+\overrightarrow{AD}=3\overrightarrow{AG}$.

A sai vì $\overrightarrow{MN}=\dfrac12\left(\overrightarrow{AD}+\overrightarrow{BC}\right)$; C sai vì $\overrightarrow{GB}=-2\overrightarrow{GN}$ chứ không phải $\overrightarrow{GD}$; D sai vì $\overrightarrow{AM}+\overrightarrow{AN}=\dfrac12\overrightarrow{AB}+\dfrac12\overrightarrow{AC}+\dfrac12\overrightarrow{AD}=\dfrac32\overrightarrow{AG}$.}
\end{ex}
